\documentclass[ ../main.tex]{subfiles}
\providecommand{\mainx}{..}
\begin{document}
\section{A theoretical \emph{cryptographic} perfect hash function}
The bit set $\{0,1\}$ is denoted by $\BitSet$.
The set of all bit strings of length $n$ is therefore $\BitSet^n$.
The cardinality of $\BitSet^n$ is $2^n$.
In the case of $\BitSet$, we denote $\BitSet^0$ by $\epsilon$.
The set of all bit strings of length $n$ or less is denoted by $\BitSet^{\leq n}$ with a cardinality $2^{n+1}-1$, e.g., $\BitSet^{\leq 2} = \epsilon \cup \BitSet \cup \BitSet^2$.
The countably infinite set of all bit strings, $\lim_{n \to \infty} \BitSet^{\leq n}$, is denoted by $\BitSet^{*}$.
A tuple $\Tuple{x_1,x_2\ldots} \in \BitSet^{*}$ is denoted a \emph{bit string}, and we typically drop the angle brackets when specifying stings, e.g., $\Tuple{x_1,x_2} \equiv x_1 x_2$.

An important operation that is closed over the free semigroup of $\BitSet$ is \emph{concatentation}, $\cat \colon \BitSet^* \times \BitSet^* \mapsto \BitSet^*$, which is defined as $a_1 \ldots a_n \cat b_1 \ldots b_m \coloneqq a_1 \ldots a_n b_1 \ldots b_m$ with special cases $x \cat \epsilon \coloneqq \epsilon \cat x \coloneqq x$.

Most hash functions are of the type $\BitSet^* \mapsto \BitSet^n$ where $n$ is some finite natural number.

Another useful function is the binary \emph{padding} function
$\Fun{pad} \colon \BitSet^* \times \IntSet_{\geq 0} \mapsto \BitSet^*$ is defined as $\Fun{pad}(x,k) \coloneqq x \cat 0^{k-\BL(x)}$ with the special case $\Fun{pad}(x,0) = \epsilon$.


The binary \emph{truncation} function $\trunc \colon \BitSet^* \times \NatSet \mapsto \BitSet^*$ is defined as $\trunc(a_1 \ldots a_k \ldots a_n, k) \coloneqq a_1 \ldots a_k$, with the special case $\trunc(x,0) \coloneqq \epsilon$.

The \emph{bit length} function $\BL \colon \BitSet^* \mapsto \NatSet$ is defined as $\BL(a_1 \ldots a_n) \coloneqq n$, e.g., if $b \in \BitSet^{n}$ then $\BL(b) = n$.

Since $\BitSet^* \cong \NatSet$, they may put into one-to-one correspondence.
A convenient one-to-one correspondence between them is given by the following definition.
\begin{definition}
	\label{def:mapping}
	The sets $\BitSet^*$ and $\NatSet$ have a bijection given by
	\begin{equation}
	b_1,\ldots,b_m \longleftrightarrow 2^m + \sum_{j=1}^{m} 2^{m - j} b_j\,.
	\end{equation}
\end{definition}

We denote the mapping described by \cref{def:mapping} with the postfix function $' \colon \BitSet^* \mapsto \NatSet$ and $' \colon \NatSet \mapsto \BitSet^*$, which are inverse functions, i.e., $x'' = x$.
An important observation of this mapping is that a natural number $n$ maps to a bit string $n'$ of length $\BL(n') = \lfloor \log_2 n \rfloor$.

More generally, if we have some set $\Set{U}$, in a physical computer there must be a way to map the elements of $\Set{U}$ to bit strings that \emph{repesent} the values.
We provide this mapping with \emph{encoder} and \emph{decoder} functions denoted respectively by $\Fun{e}[\Set{U}] \colon \Set{U} \mapsto \BitSet^*$ and $\Fun{d}[\Set{U}] \colon \BitSet^* \mapsto \Set{U}$ such that $\Fun{d}[\Set{U}] \circ \Fun{e}[\Set{U}] = \Fun{id}[\Set{U}]$ and $\Fun{e}[\Set{U}] \circ \Fun{d}[\Set{U}] = \Fun{id}[\BitSet^*]$.

A special case of the encoder and decoder is given by letting any bit string represent either the sequence of bits $a_1 \ldots a_n$ or as the binary coded decimal (BCD), e.g., $010 \longleftrightarrow 4$.
Thus, any operation on non-negative integers may be applied to bit strings without ambiguity, e.g., $010 + 01 = 11$.

Given a function of $\Fun{g} \colon \Set{X} \mapsto \Set{Y}$, the \emph{domain} of $\Fun{g}$ may be denoted by $\Dom(\Fun{g}) \coloneqq \Set{X}$ and the \emph{codomain} may be denoted by $\Codom(\Fun{g}) \coloneqq \Set{Y}$.

A random oracle as given by the following definition.
\begin{definition}
\label{def:randomoracle}
A random oracle in the family $\BitSet^{*} \mapsto \BitSet^{\infty}$ is a theoretical \emph{hash function} whose output is uniformly distributed over its codomain.
\end{definition}

The theoretical analysis of the perfect hash function makes the following assumption.
\begin{assumption}
\label{asm:ro_approx}
The hash function $\hash \colon \BitSet^{*} \mapsto \BitSet^{\infty}$ is a random oracle.
\end{assumption}

%\begin{assumption}
%\label{asm:random_sets}
%The set $\Set{S}$ consists of $m$ elements, where each element is drawn uniformly at random without replacement from $\BitSet^{*}$.
%\end{assumption}


\begin{definition}
The \emph{data type} for the \emph{cryptographic} perfect hash function under consideration is defined as $\PH \coloneqq \BitSet^* \times \NatSet$ with a value constructor $\ph \colon \PS{\Set{U}} \times [0,1] \mapsto \PH$ defined as
\begin{equation}
	\ph(\Set{X}, r) \coloneqq \Tuple{n',N}
\end{equation}
where
\begin{equation}
\begin{split}
m			&\coloneqq \Card{\Set{X}}\,,\\
N 			&\coloneqq \lceil m / r \rceil\,,\\
k			&\coloneqq \lceil\log_2 N\rceil\,,\\
\beta(x,n) 	&\coloneqq \trunc\!\left(\hash(x' \cat n'),k\right)' \mod N\,,\\
\Set{Y}[n] 	&\coloneqq \SetBuilder{\beta(x,n) \in \{0,\ldots,N-1\}}{x \in \Set{X}}\,,\\
n 			& \coloneqq \min{\SetBuilder{ j \in \NatSet }{ \Set{Y}[j] \in \PS{\NatSet} \land \Card{\Set{Y}[j]} = m}}\,.
\end{split}
\end{equation}
\end{definition}



% this is not right, N_l = N_(2^b), since if j in N_l, its bit length j' is not b, i don't think... look into it




%A family of important functions of the form $\Set{A} \times \Set{B} \mapsto \Set{A}$ and $\Set{A} \times %\Set{B} \mapsto \Set{B}$ are respectively defined as $\Fun{fst}(\Tuple{a,b}) \coloneqq a$ and %$\Fun{snd}(\Tuple{a,b}) \coloneqq b$.

Since \PH is a data type that purports to model perfect hash functions, its \emph{computational basis} is given by the following set of functions.

By the assumption of surjectivity (and the perfect hash function is not rated-distorted), then the load factor is given by $\frac{\Card{\Set{X}}}{N}$ where $N = \argmax_{x \in \Set{X}} \hash_{\Set{A}}(x)$, i.e., the maximum hash (natural ordering of integers) of $\hash_{\Set{A}}$.
\begin{definition}
The \emph{minimum} and \emph{maximum} hash of a value of type $\PH$ (that models a perfect hash function) are given respectively by $\MinHash \colon \PH \mapsto \NatSet$ and $\MaxHash \colon \PH \mapsto \NatSet$ where
\begin{equation}
	\MinHash(\Tuple{b,N}) \coloneqq 0\;\text{and}\;\MaxHash(\Tuple{b,N}) \coloneqq N-1\,.
\end{equation}
\end{definition}
%\begin{equation}
%	\MaxHash(t) \coloneqq \Fun{snd}(t)-1\,.
%\end{equation}



The most important function, the \emph{perfect hash} mapping, $\PerfHash \colon \PH \times \Set{U} \mapsto \NatSet$, is defined as
\begin{equation}
	\Fun{ph}(\Tuple{b,N},x) \coloneqq q' \mod N
\end{equation}
where
\begin{equation}
\begin{split}
	k &\coloneqq \lceil \log_2 N \rceil\,,\\
	r &\coloneqq \hash(x' \cat b)\,,\\
	q &\coloneqq \trunc(r,k)\,.
\end{split}
\end{equation}



\begin{theorem}
	% dependent-type since codomain depends on r.
A value of type $\PH$ constructed with $\ph_{\BitSet^*}(\Set{A},r)$ \emph{models} $\hash_{\Set{A}}^{r} \colon \BitSet^* \mapsto \{0,1,\ldots,k-1\}$ where $\Set{A} \subseteq \BitSet^*$ and $k = \frac{\Card{A}}{r}$.
\end{theorem}
\begin{proof}
Proof here.
\end{proof}

Suppose we have a function $\Fun{f} \colon \Set{U} \mapsto \BitSet^*$ such that $\Fun{f}\!\restriction_{\Set{A}}$ is \emph{injective}, e.g., an \emph{encoder} for values of $\Set{A} \subseteq \Set{U}$.
The composition $\hash_{\Set{A}} \coloneqq \hash_{\Set{B}} \circ \Fun{f}$ where $\Set{B} = \SetBuilder{\Fun{f}(a) \in \BitSet^*}{a \in \Set{A}}$ is a perfect hash function of type $\Set{U} \mapsto \NatSet$ over $\Set{A} \subseteq \Set{U}$.
Thus, the rest of the material in this paper does not usually assume any particular domain of the perfect hash function since injections may always be constructed for any set.

\subsection{Analysis}
Note that since $n'$ denotes the geometric code for $n$ and we choose the smallest $n$ that succeeds, where each choice of $n$ is a geometrically distributed trial with probability $p$, the expected space complexity obtains the information-theoretic lower-bound of $1.44$ bits per element.


We search \emph{all} possible hash functions which are a function of $\hash$, an approximate random oracle, and choose the \emph{perfect hash function} on a specified set that has the \emph{smallest} bit length. We describe the algorithm for performing this exhaustive search in \cref{alg:ph}.


We consider a family of perfect hash functions for a set $\Set{S}$ which are given by the output of a hash function $\hash$ that approximates a \emph{random oracle} applied to the input $x \in \Set{S}$ concatenated with a bit string $b$. We describe the generative algorithm for the perfect hash function in \cref{alg:ph}. Note that in \cref{alg:ph}, the concatenation of two bit strings $x$ and $y$ is denoted by $x \cat y$.

The perfect hash function generated by \cref{alg:ph} has the statistical property that the output is uniformly distributed as given by the following theorem.
\begin{theorem}
The perfect hash function $\hash[\Set{A}][r] \colon \Set{X} \mapsto \{0,1,\ldots,k-1\}$, where $k = \frac{\Card{\Set{A}}}{r}$, is a random oracle over $\SetDiff[\Set{X}][\Set{A}]$ and the restriction $\hash[\Set{A}][r] \!\restriction_{\Set{A}} \colon \Set{A} \mapsto \{0,1,\ldots,k-1\}$ is a random $k$-permutation oracle of $\Set{A}$.
\end{theorem}
\begin{proof}
First, in \cref{alg:ph} on Line~5, $b_n$ is a bit string such that each $x \in \Set{S}$ concatenated with $b_n$ hashes to a unique integer in $\{1,\ldots,N\}$ by the hash function $\hash$, thus the hash function found perfectly hashes the elements of $\Set{S}$. 

Second, by \cref{def:randomoracle}, $\hash$ approximates a random oracle whose output is uniformly distributed over the elements of $\BitSet^{n}$. Viewing the elements of $\BitSet^{n}$ as integers, where $\Card{\BitSet^{n}} = 2^n$, $\hash$ is uniformly distributed over $\{0,1,\ldots,2^n-1\}$.

If $N = k 2^n$ for some integer $k$, the remainder of the output of $\hash$ when dividing by $N$, given by the modulo operator, and adding $1$ is uniformly distributed over $\{1,\ldots,N\}$ since each $j \in \{1,\ldots,N\}$ has an equal number of hashes assigned to it by $\hash$. If $2^n \gg N$ and $N \neq k 2^n$ for some integer $k$, then it is approximately uniformly distributed and converges to the uniform distribution as $n \to \infty$.
\end{proof}

The probability that no collisions occur for a particular bit string in \cref{alg:ph} is given by the following theorem.
\begin{theorem}
The probability that a bit string $b \in \BitSet^{*}$ results in a perfect hash function is given by
\begin{equation}
\label{eq:p_m_N}
	p(m,r) = N^{-m} P^{N}_{m}
%    p(m,N) = \frac{N!}{N^m (N-m)!}\,,
\end{equation}
where $m$ is the cardinality of the set being perfectly hashed, $r$ is the load factor, and $N \coloneqq \frac{m}{r}$.
\end{theorem}
\begin{proof}
Suppose we have a set $\Set{A} \subset \BitSet^*$ of cardinality $m$.
The set of perfect hash functions $\hash_{\Set{A}}^{r} \colon \BitSet^* \mapsto \{0,1,\ldots,N-1\}$ restricted to $\Set{A}$ where $N \coloneqq \frac{m}{r}$ has a cardinality given by $P^{N}_{m}$ since any choice of $m$ out of $N$ elements in the codomain and any ordering of the $m$ elements in $\Set{A}$ to the chosen $m$ elements in $\{0,1,\ldots,N-1\}$ satisfy the definition of perfect hashing.

%First, we must choose $m$ out of the $N$ elements in the codomain to perfectly hash the $m$ elements in $\Set{A}$.
%There are a total of $N \choose m$ ways of doing this.
%Once we have chosen the $m$ elements, we do not care about the order in which we assign the $m$ elements in $\Set{A}$ %to the $m$ elements in the codomain.
%There are a total of $m!$ such orderings, any of which satisfy the perfect hashing criteria.
%Thus, by the product rule, there are ${N \choose m} m! = P_{m}^{N}$ ways to make these choices.

The set of hash functions $\BitSet^* \mapsto \{0,1,\ldots,N-1\}$ restricted to $\Set{A}$ has a cardinality of $N^m$.
Therefore, the ratio of perfect hash functions restricted to $\Set{A}$ to the total hash functions restricted to $\Set{A}$ is just
\begin{equation}
	p = \frac{P^{N}_{m}}{N^m}\,.
\end{equation}

By the property of the hash function $\hash^*$ being a random oracle, the algorithm defined in \cref{dummy} randomly samples one of the hash functions, which is a perfect hash function with probabilty $p$.

%$\hash_{\Set{A}}^{r}$ to $\Set{A}$, there are $N^m$ hash functions of this type. type $\BitSet^* \mapsto \NatSet$ over $\Set{A}$ with a load factor $r$ restricted to $\Set{A}$ with 
	
%Suppose we have a set $\Set{S} = \{x_1,\ldots,x_m\}$. First, we consider a singular set consisting of $x_1$. The probability that no collision occurs is $1$ since there are no other members to collide with. Second, we add $x_2$ to the set. The probability that no collision occurs is given by the probability
%\begin{equation}
%    \Prob{\hash(x_2 \cat b) \neq \hash(x_1 \cat b)} = \frac{N-1}{N}\,.
%\end{equation}
%This probability follows from the fact that out of the $N$ hashes to choose from, $N-1$ result in no collision. %Extending this logic, the probability that no collision occurs for all $m$ elements is given by
%\begin{align}
%    p(m,N)
%        &= \prod_{i=0}^{m-1} \frac{N - i}{N}\\
%        &= \frac{N!}{N^m (N-m)!}\,.
%\end{align}
\end{proof}

\Cref{eq:p_m_N} may be reparameterized with respect to the load factor. By \cref{eq:loadfactor}, $r = m/N$. Solving this equation with respect to $N$ yields the solution
\begin{equation}
    N = \frac{m}{r}\,,
\end{equation}
and thus plugging in this value of $N$ yields the result
\begin{equation}
\label{eq:p_m_r}
    p(m,r) = \frac{\left(\frac{m}{r}\right)!}{\left(\frac{m}{r}\right)^m \left(\frac{m}{r}-m\right)!}\,.
\end{equation}

The expected \emph{in-place} coding size is given by the following theorem.
\begin{theorem}
The expected coding size is given approximately by
\begin{equation}
\label{eq:code_size}
    \log_2 \mathrm{e} - \left(\frac{1}{r} - 1\right) \log_2\!\left(\frac{1}{1-r}\right) \si{bits \per element}\,.
\end{equation}
\end{theorem}
\begin{proof}
\begin{equation}
    \frac{n}{r}(1 - r) \log_2 \left(1 - r\right) - \log_2 n - (u - n) \log_2 (1 - n/u)
\end{equation}
\begin{equation}
    \left(\frac{1}{r} - 1\right) \log_2 \left(1 - r\right) - \frac{1}{n}\log_2 n - \frac{u - n}{n} \log_2\left(1 - \frac{n}{u}\right)
\end{equation}

The space required for the perfect hash function found by \cref{alg:ph} is of the order of the length $n$ of the bit string $b_n$ in the returned tuple. Therefore, for space efficiency, the algorithm exhaustively searches for a bit string in the order of increasing size $n$.

We are interested in the first case when no collision occurs, which is a geometric distribution with probability of success $p(m,r)$ as given by the discrete random variable
\begin{equation}
    \RV{Q} \sim \geodist\!\bigl(p\!\left(m,r\right)\bigr)\,,
\end{equation}
where $p(m,r)$ is given by \cref{eq:p_m_r}. The expected number of trials for the geometric distribution is given by
\begin{equation}
    \Expect{\RV{Q}} = \frac{1}{p(m,r)} = \frac{\left(\frac{m}{r}\right)^m \left(\frac{m}{r}-m\right)!}{\left(\frac{m}{r}\right)!}\,.
\end{equation}
By \cref{def:mapping}, the \nth trial uniquely maps to a bit string of length $m = \lfloor \log_2 n \rfloor$. Thus, the expected bit length is given approximately by
\begin{align}
    \Expect{\log_2 \RV{Q}}
        &= \log_2\!\left(\frac{\left(\frac{m}{r}\right)^m \left(\frac{m}{r}-m\right)!}{\left(\frac{m}{r}\right)!}\right)\\
        &= m \log_2\!\left(\frac{m}{r}\right) +
            \log_2\!\left(\frac{m}{r}-m\right)! - \log_2\!\left(\frac{m}{r}\right)! \; \si{bits}\,.
\end{align}
By Stirling's approximation, 
\begin{equation}
    \log_2 n! \approx n \left(\log_2 n - \log_2 \mathrm{e}\right)\,.
\end{equation}
and so the expectation may be rewritten approximately as
\begin{equation}
\begin{split}
    \Expect{\RV{Q}} \approx &\;
        m \log_2\!\left(\frac{m}{r}\right) -
            \frac{m}{r} \left(\log_2\!\left(\frac{m}{r}\right) - \log_2 \mathrm{e}\right) +\\
        &\qquad\left(\frac{m}{r}-m\right) \left(\log_2\!\left(\frac{m}{r}-m\right) - \log_2 \mathrm{e} \right) \; \si{bits}\,.
\end{split}
\end{equation}
Since we are interested in the expected \emph{bits per element}, we divide the expectation by $m$ and after further simplification arrive at
\begin{equation}
\begin{split}
    &\log_2\!\left(\frac{m}{r}\right) - \frac{1}{r} \log_2\!\left(\frac{m}{r}\right) +\\
    &\qquad \left(\frac{1}{r}-1\right) \log_2\!\left(\frac{m}{r}-m\right) + \log_2 \mathrm{e} \; \si{bits \per element}\,.
\end{split}
\end{equation}
After further simplification, we arrive at
\begin{equation}
    \log_2 \mathrm{e} - \left(\frac{1}{r} - 1\right) \log_2\!\left(\frac{1}{1-r}\right) \si{bits \per element}\,.
\end{equation}
\end{proof}

The entropy of $\RV{Q}$ is given by
\begin{equation}
    \frac{-(1-p)}{p} \log_2(1 - p) - \log_2(p)
\end{equation}

\begin{equation}
    \left(\frac{1}{p}-1\right) \log_2\!\left(\frac{1}{1 - p}\right) + \log_2\frac{1}{p}
\end{equation}


Note that according to \cref{??}, \cref{alg:ph} has an exponential time 
complexity with respect to the cardinality of the input set $\Set{S}$. As a 
result, \cref{alg:ph} is not a practical algorithm for any reasonably large 
$m$. However, it is intended to illustrate theoretical properties useful to 
data structures that implement \emph{oblivious} sets and maps\cite{phf,pmf} 
based on the perfect hash function. Simple and efficient algorithms exist 
\cite{chd,simple}. 

The theoretical lower-bound of a \emph{minimal} perfect hash function is given by the following postulate.
\begin{postulate}
\label{pst:lb}
The theoretical lower-bound for \emph{minimal} perfect hash functions has an expected coding size given approximately by
\begin{equation}
\label{eq:low_bound}
    \SI{1.44}{bits \per element}\,.
\end{equation}
\end{postulate}

\begin{theorem}
The cryptographic perfect hash generator given by \cref{alg:ph} results in a \emph{minimal} perfect hash function that obtains the theoretical lower-bound of $\SI{1.44}{bits \per element}$ by invoking \makehash{$\cdot$, $r=1$}.
\end{theorem}
\begin{proof}
By \cref{eq:code_size}, letting $r \to 1^{+}$ for the \emph{minimal} perfect hash results in an expected coding size given by
\begin{equation}
    \log_2 \mathrm{e} - \lim_{a \to 0^{-}} a \log_2 \frac{1}{a} = 1.44 \; \si{bits \per element}\,,
\end{equation}
which is the expected lower-bound given by \cref{pst:lb}.
\end{proof}
This is as expected, since \cref{alg:ph} finds the \emph{smallest} perfect hash function that is a function of a random oracle for any load factor $0 < r \leq 1$ in which the distribution of the sets are uniformly distributed. The following corollary follows as a result.
\begin{corollary}
The lower-bound for perfect hash functions with a load factor $0 < r \leq 1$ is given by \cref{eq:code_size}.
\end{corollary}
As $r$ goes to $0$, the expected bits per element goes to $0$.
t
Note that the \emph{probability mass function} of the random bit string found for $b$ is known exactly. Thus, if the desire is to \emph{serialize} the perfect hash function for transmission or storage, shorter bit strings may be assigned to more probable bit strings. This representation is not usable \emph{in-place}, i.e., the serialization must be decoded, but it can reduce transmission or storage cost.

\subsection{Entropy}

The entropy of the random bit length $\RV{N}$ is given by
\begin{equation}
?    
\end{equation}
\begin{proof}
\begin{align}
    -\sum_{n=0}^{\infty} \log_2\!\left(\pmf_{\RV{N}}(n \Given m,r)\right) \pmf_{\RV{N}}(n \Given m,r)\,.
\end{align}
\end{proof}

\end{document}